A. QUALITATIVE JUDGMENT

The immediate judgment is the judgment of existence:
the subject posited in a universality, as its predicate,
which is an immediate (thus sensory) quality.
(1) Positive judgment: the individual is a particular.
But the individual is not a particular;
more precisely, such an individual quality
does not correspond to the concrete nature of the subject;
(2) negative judgment.

    It is one of the most essential logical prejudices that such qualitative
    judgments as 'the rose is red' or 'the rose is not red' can contain truth.
    They can be correct, i.e. in the limited sphere of perception,
    finite representing, and thinking. This depends upon the content,
    which is just as much a finite content, untrue for itself.
    But the truth rests solely on the form, i.e. the posited concept and
    the reality corresponding to it; but such truth is not at hand
    in the qualitative judgment.

Addition. Correctness and truth are very frequently considered to mean the same
thing in ordinary life and one accordingly speaks of the truth of some content
where it is a matter of mere correctness. Correctness generally affects merely the
formal agreement of our representation with its content; however this content may
be otherwise constituted. The truth consists, by contrast, in the agreement of the
object with itself, i.e. with its concept. It may be correct anyway that someone is
sick or that someone has stolen something. But such content is not true since a
sick body is not in agreement with the concept of life, and so too theft is an action
that does not correspond to the concept of human action. What is to be taken
from these examples is that an immediate judgment, in which an abstract quality
is asserted of something immediately individual, however correct it might be, stiII
contains no truth since subject and predicate do not stand in the judgment in
the connection of reality and concept to one another. The lack of truth of the
immediate judgment consists, further, in the fact that its form and content do not
correspond to one another. If we say 'this rose is red', then it lies in the copula 'is'
that subject and predicate agree with one another. But now the rose, as something
concrete, is not merely red; instead it also has an odour, a determinate form,
and many other sorts of determinations that are not contained in the predicate
'red'. On the other side, this predicate, as an abstract universal, does not apply
merely to this subject. There are also, in addition, other flowers and generally other
objects that are likewise red. Subject and predicate in the immediate judgment
thus come into contact with one another, as it were, only at one point but.They
do not cover one another. The state of affairs is quite different in the conceptual
judgment. If we say 'this action is good', this is then a conceptual judgment.
One notices immediately that here, between subject and predicate, there is not
this loose and external connection as there is in the immediate judgment. In the
immediate judgment the predicate consists in some abstract quality or other which
mayor may not apply to the subject. In the conceptual judgment, by contrast, the
predicate is, as it were, the soul of the subject, by means of which the subject, as
the body of this soul, is determined through and through.

In this as first negation there still remains
the relation of the subject to the predicate,
which is thereby something relatively universal,
the determinacy of which has only been negated
('the rose is not red' entails that it still has colour;
immediately another [colour] which, however,
would only be a positive judgment in turn).
The individual, however, is also not a universal.
(3) Hence, (aa) the judgment collapses in itself
into the empty identical relation:
the individual is the individual - identical judgment;
and (bb) it collapses into itself as the present,
complete inadequacy of the subject and predicate:
a so-called infinite judgment.

    Examples of the latter are 'the spirit is no elephant',
    'a lion is no table', and so forth - sentences that are correct
    but as nonsensical as the identical sentences 'a lion is a lion',
    'the spirit is spirit'. These sentences are, to be sure,
    the truth of the immediate, so-called qualitative judgment,
    but not judgments at all, and they can only surface in
    a subjective thinking that can fix upon an untrue abstraction.
    Objectively considered, they express the nature of beings
    or sensory things, namely, that they collapse into an
    empty identity and into a fulfilled relation that is, however,
    the qualitative otherness of what is related, their complete inadequacy.

Addition. The negative-infinite judgment, in which no relation at all between
subject and predicate is on hand any more, is usually cited in formal logic merely
as a senseless curiosity. Nevertheless, this infinite judgment is in fact not
to be considered merely as a contingent form of subjective thinking. Instead it
ensues as the very next dialectical result of the preceding, immediate judgments (of
the positive and the simply negative), whose finitude and lack of truth explicitly
come to light in it. Crime can be regarded as an objective example of the negative-
infinite judgment. Whoever commits a crime, more precisely a theft, does not
merely negate, as in the civil juridical dispute, the particular right of someone
else to this specific matter. Instead he negates the right of that person altogether
and, for this reason, he is not merely ordered to restore the matter which he stole,
but is instead punished in addition because he violated the right as such, i.e.
the right in general. The civil juridical dispute is, by contrast, an example of the
simple-negative judgment since in it merely this particular right is negated and
right in general is recognized in the process. The connection here is thus as it is
for the negative judgment 'this flower is not red', by means of which merely this
particular colour, but not colour altogether, is negated in regard to the flower since
it can still be blue, yellow, and so forth. Likewise then, too, death is a negative-
infinite judgment in contrast to sickness, which is a simple-negative judgment. In
a sickness, merely this or that particular vital function is restricted or negated; by
contrast, in death, as one would say, body and soul separate from one another, i.e.
subject and predicate fall completely outside one another.

B. THE JUDGMENT OF REFLECTION

The individual, posited as individual (reflected in itself) in the judgment,
has a predicate, opposite which the subject, relating itself to itself,
remains at the same time an other.
In the concrete existence the subject is no longer immediately qualitative,
but is instead in a connection with and joined to an other, an external world.
The universality has acquired hereby the meaning of this relativity.
(For example, useful, dangerous; weight, acidity, - then drive, and so forth.)

Addition. The judgment of reflection is distinguished generally from the
qualitative judgment by the fact that its predicate is no longer an immediate, abstract
quality but instead of the sort that, by means of it, the subject demonstrates itself
to be related to the other. If we say, for example, 'this rose is red', we consider the
subject in its immediate individuality without relation to another. If, by contrast,
we make the judgment 'this plant has healing powers', we consider the subject, the
plant, as standing in relation with another (the illness to be healed by it) through
its predicate, the healing capacity. Matters are similar with the judgments 'this
body is elastic', 'this instrument is useful', 'this punishment works as a deterrent' ,
and so forth. The predicates of such judgments are generally determinations of
reflection, by means of which one has moved beyond the immediate individuality
of the subject while the concept of it is still not given.
The usual sort of rationalizing tends above all to immerse itself in this manner of judgment.
The more concrete the object of concern, the more viewpoints it presents for
reflection, by means of which, meanwhile, the distinctive nature, i.e. its concept,
is not exhausted.

(1) The subject, the individual as individual (in the singular judgment), is a universal.
(2) In this relation it is elevated above its singularity.
This expansion is an external one, the subjective reflection,
at first the indeterminate particularity
(in the particular judgment which is, immediately, negative as well as positive;
the individual is in itself divided, it relates itself in part to itself, in part to another).
(3) Some are the universal, so the particularity is expanded to universality;
or this universality, determined by the individuality of the subject,
is the set of all (commonality, the usual universality-of-reflection).

Addition. Since it is determined in the singular judgment as universality, the
subject by this means moves beyond itself, past itself as this mere individual.
When we say 'this plant has healing powers', this entails not merely that this
individual plant has them but that several or some do and this results in the
particular judgment (,Some plants have healing powers', 'Some human beings are
inventive', and so forth). Through this particularity, the immediately individual
[subject] loses its self-sufficiency and enters into a connection
with another. The human being is, as this human being, no longer merely this
individual human being; instead he stands alongside other human beings and is
thus one in a group. Precisely by this means, however, it also belongs to
the universal and is thereby elevated. The particular judgment is positive as well
as negative. If only some bodies are elastic, then the rest are not elastic. - Herein
lies, then, again the progression to the third form of the judgment-of-reflection,
i.e. to the judgment of the set of all ('all human beings are mortal', 'all metals are
conductors of electricity'). The set of all is that very form of universality
towards which reflection at first tends. In this connection the individuals form
the foundation and it is our subjective act by means of which the individuals
are gathered together and are determined [as belonging together] in their entirety.
The universal appears here only as an external bond which
encompasses the individuals subsisting for themselves and indifferent to it. The
universal is, nevertheless, in fact the ground and basis, the root and substance of the
individual. If we consider, for example, Caius, Titus, Sempronius, and the other
inhabitants of a city or a country, then the fact that they are collectively human
beings is not merely something common to them, but their universal their genus,
and all these individuals would not be at all without this, their genus. In contrast
to this, matters are different with that superficial, only so-called universality that
is nn fact something that merely accrues to all individuals and is common to them.
It has been noted that human beings, in contrast to animals, have this in common
with one another, that they are equipped with ear lobes. It is, meanwhile, apparent
that if, somehow, one or the other should not have ear lobes, the rest of his being,
hiS character, his capacities, and so forth would not be affected by this. It would,
by contrast, make no sense to assume that Caius could somehow not be a human
being but be brave, learned, and so forth. What the individual human being is in
particular, this is only insofar as he is, above all, a human being as such and in the
universal sense, and this universal is not only something external
to and alongside other abstract qualities or mere determinations of reflection.
Instead it is much more what pervades everything particular, encompassing it
within itself.

By the fact that the subject is likewise determined as universal,
the identity of it and the predicate is posited as indifferent,
as is, thanks to this, the determination of the judgment itself.
This unity of the content as the universal identical with
the subject's negative reflection-in-itself makes
the relation of the judgment a necessary relation.

Addition. The progression from the reflexive judgment of the set of all to the
necessary judgment can be found already in our ordinary consciousness insofar as we say:
'what accrues to everything, accrues to the genus and is, therefore, necessary.'
When we say: 'all plants', 'all humans', and so forth, this is the same as
if we say 'the plant', 'the human', and so forth.

C. JUDGMENT OF NECESSITY

The judgment of necessity as the identity of the content in its difference
(1) contains within the predicate in part the substance or nature of the
subject, the concrete universal - the genus; in part, since this universal
equally contains in itself the determinacy as negative, the excluding essential
determinacy - the species; - categorical judgment.
(2) In keeping with their substantiality, the two sides acquire the form
of self-sufficient actuality, the identity of which is only an inner identity,
and with that the actuality of the one is at the same time not its actuality,
but instead the being of the other; ~ hypothetical judgment.
(3) At the same time, in this externalization of the concept, the inner
identity is posited and so the universal is the genus that is identical with
itself in its excluding individuality. The judgment which has this universal
on both sides of it, the one time as such, the other time as the sphere of
its self-excluding particularization - the either/or of which just as much
as the as well as is the genus - is the disjunctive judgment. With this, the
universality at first as genus and then also as the scope of its species is
determined and posited as a totality.

Addition. The categorical judgment ('Gold is a metal', 'The rose is a plant') is
the immediate judgment of necessity and corresponds, in the sphere of essence, to
the relationship of substantiality. All things are a categorical judgment, i.e. they
have their substantial nature, which forms the fixed and unchangeable foundation
of them. Only when we regard things from the viewpoint of their genus and as
determined by it with necessity, does the judgment begin to be a true one. It must
be designated a deficiency in someone's training in logic, if judgments like these:
'Gold is expensive' and 'gold is a metal' are regarded as standing on the same level.
That gold is expensive concerns an external relation of it to our inclinations and
needs, to the costs of acquiring it, and so on, and the gold remains what it is, even
if that external relation alters or falls away. By contrast, being a metal constitutes
the substantial nature of gold, without which it or anything else that is otherwise
in it or asserted of it cannot subsist. Matters are the same if we say 'Caius is a
human being'; in this way we declare that everything that he may otherwise be
only has value and meaning insofar as it corresponds to this, his substantial nature,
to be a human being. - Furthermore, however, even the categorical judgment
remains deficient insofar as in it the factor of particularity does not yet receive
its due. Thus, for example, the gold is indeed metal, but silver, iron, and
so forth are likewise metals, and being metal as such behaves indifferently to the
particular character of its species. Herein lies the progression from the categorical
to the hypothetical judgment which can be expressed by the formula: 'if A is,
then B is'. We have here the same progression as earlier from the relationship
of substantiality to the relationship of causality. In the hypothetical judgment,
the determinacy of the content appears as mediated, as dependent upon another,
and this is then precisely the relationship of cause and effect. The meaning of
the hypothetical judgment is then generally this, that through it the universal is
posited in its particularization and, with this, we acquire, as the third form of
necessary judgment, the disjunctive judgment. 'A is either B or C or D'; the poetic
artwork is either epic or lyrical or dramatic; the colour is either yellow or blue or
red, and so on. The two sides of the disjunctive judgment are identical. The genus
is the totality of its species and the totality of the species is the genus. This unity
of the universal and the particular is the concept and it is this, which now forms
the content of the judgment. of its existence, the subject then expresses the relation of
that particularity to its constitution, I.e. to its genus and, with this, expresses
what (see preceding section) makes up the content of the predicate (this-
the immediate individuality - house -'genus -, so and so constituted - partic-
ularity -, is good or bad) - apodictic judgment. - All things are a genus (their
determination and purpose) in one individual actuality with a particuktr
constitution; and their [i.e. all things'] finitude consists in the fact that
their particular [character] mayor may not be adequate to the universal

D. THE JUDGMENT OF THE CONCEPT

The judgment of the concept has the concept, the totality in simple form,
for its content, the universal with its complete determinacy.
The subject is (1) initially an individual that has, as its predicate,
the reflection of the particular existence on its universal,
the agreement or lack of agreement of these two determinations:
good, true, correct, and so forth - assertoric judgment.

    In ordinary life, too, one only calls it judging when a judgment is of
    this sort, e.g. the judgment whether an object, action, and so forth is
    good or bad, true, beautiful, and so forth. One will not ascribe a
    power of judgment to someone [simply] for knowing, for example,
    how to make positive or negative judgments such as 'this rose is red',
    'this painting is red, green, dusty', and so forth.
    Even in philosophy, through the principle of immediate knowing
    and believing, the assertoric judgment has been made into the sole
    and essential form of the doctrine (despite the fact that in society the
    assertoric judgment counts as improper, when someone claims that
    it is supposed to be valid by itself). In the so-called philosophical
    works that maintain that principle, one can read hundreds upon
    hundreds of assurances about reason, knowing, thinking,
    and so forth, which seek to gain credence for themselves through
    endless repetitions of one and the same point, since external
    authority no longer counts for much.

In what is at first the immediate subject of the assertoric judgment,
this judgment does not contain that relation of the particular and the universal.
that is expressed in the predicate.
This judgment is thus merely a subjective particularity and
the opposite assurance stands over against it with the
same right or, rather, the same lack of right.
It is thus (2) at the same time only a problematic judgment.
But (3) [insofar as] the objective particularity is posited in the subject,
its particularity as the constitution of its existence,
the subject then expresses the relation of that particularity to its constitution,
i.e. to its genus and, with this, expresses what (see preceding section)
makes up the content of the predicate
(this - the immediate individuality - house - genus - so and so constituted - particularity,
is good or bad) - apodictic judgment.
All things are a genus (their determination and purpose)
in one individual actuality with a particular constitution;
and their [i.e. all things'] finitude consists in the fact that
their particular [character] may or may not be adequate to the universal.
